\chapter{Materiales y métodos}
\label{cap:metodos}

\section{Visión general del \emph{pipeline}}
\label{sec:pipeline_overview}

El framework \framework\ se organiza como una tubería de pasos secuenciales
que producen artefactos intermedios (matrices, CSV, JSON) versionados en
disco. Cada paso puede ejecutarse de forma aislada o encadenado desde el
orquestador. La figura conceptual del pipeline es la siguiente:

\begin{enumerate}
  \item Descarga de cohortes GEO y mapeo sonda $\to$ símbolo génico.
  \item Normalización por estudio (log$_2$ y cuantiles).
  \item Curación clínica con LLM (Llama~3.3-70b).
  \item Aplicación de criterios de inclusión y auditoría.
  \item Análisis diferencial (Welch + FDR).
  \item Clasificación supervisada (LASSO L1, RF, SVM).
  \item Validación externa LODO.
  \item Meta-análisis por consenso y firma génica validada.
  \item Panel mínimo y validación contra IHC clínica.
  \item Interfaz web y asistente conversacional.
\end{enumerate}

La totalidad del \emph{pipeline} está implementada en Python 3.12 y hace
uso de las librerías científicas estándar (\texttt{pandas}, \texttt{numpy},
\texttt{scikit-learn}, \texttt{scipy}, \texttt{statsmodels}). La
reproducibilidad se garantiza mediante fijación de la semilla aleatoria
(\texttt{random\_state=42}) en todos los pasos con componente estocástico
y mediante el registro de cada resultado intermedio en disco. El código
fuente se aloja en un repositorio Git público
(\url{https://github.com/danitapiadiez-gif/TFM-Bioinformatica}).

La figura~\ref{fig:pipeline} representa gráficamente la arquitectura del
pipeline, agrupando los pasos por fase y señalando los artefactos
principales que fluyen entre ellos.

\begin{figure}[H]
  \centering
  \begin{tikzpicture}[node distance=0.55cm and 0.9cm]
    % Fila 1: entrada
    \node[entrada] (geo)      {NCBI GEO};
    \node[paso, right=of geo] (descarga) {\textbf{1} Descarga};
    \node[paso, right=of descarga] (norma) {\textbf{2} Normalización};
    % Fila 2: curación
    \node[paso, below=of descarga] (llm)  {\textbf{3} Curación LLM\\(Llama 3.3-70b)};
    \node[paso, right=of llm]     (crit) {\textbf{4} Criterios de\\inclusión};
    % Fila 3: analisis
    \node[paso, below=of llm]     (dif)  {\textbf{5} Diferencial\\Welch + FDR};
    \node[paso, right=of dif]     (ml)   {\textbf{6} Clasificación\\LASSO L1 · RF · SVM};
    % Fila 4: LODO
    \node[paso, below=of dif]     (lodo) {\textbf{7} LODO\\externa};
    \node[paso, right=of lodo]    (cons) {\textbf{8} Consenso\\multi-cohorte};
    % Fila 5: entregables
    \node[salida, below=of lodo]  (firma) {\textbf{9} Firma\\1174 genes};
    \node[salida, right=of firma] (panel) {\textbf{9} Panel\\mínimo 20};
    \node[salida, right=of panel] (ihc)   {\textbf{9} Validación\\IHC 18/20};

    % Flechas principales
    \draw[flecha] (geo) -- (descarga);
    \draw[flecha] (descarga) -- (norma);
    \draw[flecha] (norma.south) |- (llm.east);
    \draw[flecha] (llm) -- (crit);
    \draw[flecha] (crit.south) |- (dif.east);
    \draw[flecha] (dif) -- (ml);
    \draw[flecha] (ml.south) |- (lodo.east);
    \draw[flecha] (lodo) -- (cons);
    \draw[flecha] (cons.south) |- (firma.east);
    \draw[flecha] (firma) -- (panel);
    \draw[flecha] (panel) -- (ihc);
  \end{tikzpicture}
  \caption{Arquitectura del \emph{pipeline} de \framework. Los pasos con
  fondo verde son entradas (repositorios públicos), los pasos naranjas
  son entregables (firma, panel mínimo y validación externa contra IHC
  clínica), y los pasos azules son transformaciones intermedias.}
  \label{fig:pipeline}
\end{figure}

\section{Adquisición y normalización de cohortes}
\label{sec:adquisicion}

\subsection{Selección de cohortes}
\label{sec:seleccion_cohortes}

Se seleccionaron 11 cohortes del repositorio NCBI GEO que cumplían los
siguientes requisitos previos: (i) especie \emph{Homo sapiens}; (ii) tejido
pulmón; (iii) plataforma Affymetrix HG-U133 Plus 2.0 (GPL570) o compatible;
(iv) tamaño muestral mínimo de 10 muestras; (v) disponibilidad pública
completa (no restricted access). La tabla~\ref{tab:cohortes} enumera las
cohortes seleccionadas.

\begin{table}[H]
  \centering
  \small
  \caption{Cohortes GEO utilizadas en el trabajo.}
  \label{tab:cohortes}
  \begin{tabular}{lrrrrl}
    \toprule
    \textbf{Cohorte} & \textbf{Muestras totales} & \textbf{Sanas} & \textbf{Enfermas} & \textbf{Sin clasificar} & \textbf{Evaluable} \\
    \midrule
    GSE118370 & 12 & 6 & 6 & 0 & sí \\
    GSE140797 & 14 & 0 & 7 & 7 & no \\
    GSE18842  & 91 & 45 & 46 & 0 & sí \\
    GSE19188  & 156 & 65 & 91 & 0 & sí \\
    GSE19804  & 120 & 60 & 60 & 0 & sí \\
    GSE23066  & 10 & 5 & 5 & 0 & sí \\
    GSE30219  & 307 & 0 & 307 & 0 & no \\
    GSE31210  & 246 & 20 & 226 & 0 & sí \\
    GSE40791  & 194 & 10 & 2 & 182 & sí \\
    GSE50081  & 181 & 0 & 181 & 0 & no \\
    GSE7670   & 15 & 8 & 7 & 0 & sí \\
    \midrule
    \textbf{Total} & \textbf{1\,346} & \textbf{219} & \textbf{938} & \textbf{189} & \textbf{8 / 11} \\
    \bottomrule
  \end{tabular}
\end{table}

\subsection{Descarga y mapeo}
\label{sec:descarga}

La descarga se realiza mediante la librería \texttt{GEOparse}, que consulta
la API de GEO y descarga tanto la matriz SOFT como los ficheros
suplementarios cuando existen. Para cada cohorte se obtiene:

\begin{itemize}
  \item La matriz de expresión sonda~$\times$~muestra.
  \item Los metadatos de muestra (\texttt{characteristics\_ch1},
  \texttt{title}, \texttt{source\_name}), preservados como texto libre.
  \item La tabla de anotación de la plataforma, con la correspondencia
  sonda~$\to$~símbolo génico oficial (HUGO).
\end{itemize}

El mapeo sonda~$\to$~gen agrega las sondas múltiples que apuntan al mismo
símbolo tomando el promedio geométrico de sus intensidades. Las sondas sin
anotación se descartan.

\subsection{Normalización}
\label{sec:normalizacion}

La normalización se realiza \emph{dentro de cada estudio}, sin intentar
corrección de efecto lote entre cohortes. Los pasos son:

\begin{enumerate}
  \item Transformación logarítmica: $x \gets \log_2(x + 1)$.
  \item Normalización por cuantiles: cada muestra se ajusta al cuantil
  correspondiente de la distribución media de la cohorte.
\end{enumerate}

La razón de no aplicar corrección conjunta de lote (por ejemplo ComBat) es
que LODO no requiere que el efecto lote esté corregido: lo mide. Si la
firma resulta robusta bajo LODO, es porque los genes seleccionados
mantienen su patrón de cambio a pesar del efecto lote, lo que es una
propiedad deseable. Aplicar ComBat sobre cohortes con distinta prevalencia
de la variable de interés puede introducir sesgos difíciles de auditar.

\section{Curación clínica con modelo de lenguaje}
\label{sec:curacion_llm}

\subsection{Formulación de la tarea}
\label{sec:tarea_llm}

Para cada muestra se construye un \emph{prompt} que incluye:

\begin{itemize}
  \item El texto libre de \texttt{title}, \texttt{source\_name} y todas las
  \texttt{characteristics\_ch1} de la muestra.
  \item Una instrucción canónica que solicita la clasificación en
  \texttt{sano}, \texttt{enfermo} o \texttt{ambiguo}, con criterios
  explícitos y ejemplos de cada categoría.
  \item Una restricción de formato de salida: JSON con dos campos
  (\texttt{etiqueta} y \texttt{motivo}).
\end{itemize}

La categoría \texttt{ambiguo} recoge tanto muestras cuya descripción no
permite asignar grupo (por ejemplo, línea celular sin especificar
condición) como muestras que pertenecen a subtipos no contemplados en el
estudio (metástasis, tumores neuroendocrinos si la cohorte se enfoca en
NSCLC, etc.). Este diseño evita forzar clasificaciones erróneas y produce
un descarte explícito.

\subsection{Modelo y servicio}
\label{sec:modelo_llm}

Se utiliza Llama~3.3-70b servido a través de la API pública de Groq. Las
peticiones se envían con \texttt{temperature=0.0} para maximizar el
determinismo. La latencia media por muestra es del orden de 400\,ms y el
coste marginal es inferior a 0,001~USD por muestra al precio actual de la
tier gratuita del servicio.

\subsection{Auditoría de la curación}
\label{sec:auditoria_curacion}

Para cada cohorte se registran las siguientes métricas de curación en el
fichero \texttt{AUDITORIA\_COHORTES.csv}:

\begin{itemize}
  \item Número total de muestras descargadas.
  \item Muestras clasificadas como \texttt{sano}, \texttt{enfermo},
  \texttt{sin\_clasificar}.
  \item Tasa de éxito de curación: proporción de muestras con etiqueta
  distinta de \texttt{ambiguo}.
  \item Marca \texttt{Evaluable\_Como\_Test}: booleano que indica si la
  cohorte tiene al menos una muestra de cada clase (sano y enfermo).
\end{itemize}

\section{Criterios de inclusión}
\label{sec:criterios_inclusion}

Una cohorte se considera \emph{evaluable} para la tarea tumor \emph{vs.}
sano si cumple simultáneamente:

\begin{enumerate}
  \item Presencia de al menos una muestra sana y una enferma.
  \item Tasa de éxito de curación superior al 40\,\% (por debajo de este
  umbral, la señal está dominada por descartes y no permite entrenar de
  forma fiable).
  \item Alineamiento verificado entre matriz de expresión y metadatos por
  identificador \texttt{geo\_accession}. Cualquier discordancia se
  registra en \texttt{N\_Muestras\_Desalineadas} y bloquea la cohorte.
\end{enumerate}

Sobre las 11 cohortes descargadas, 8 cumplen los tres criterios. Las tres
excluidas (GSE30219, GSE50081, GSE140797) no tienen controles sanos
evaluables. Sus muestras tumorales sí se aprovechan en el entrenamiento
LODO (aportan variabilidad a la clase mayoritaria), pero no pueden actuar
como test.

Para la tarea de subtipo histológico ADC \emph{vs.} SQC, el criterio de
evaluabilidad se reformula como: presencia de al menos 10 muestras de cada
subtipo, sobre plataforma GPL570 (para maximizar la comparabilidad técnica).
Este criterio lo cumplen 3 cohortes: GSE19188, GSE30219 y GSE50081, con
un total de 388 muestras (146+170+72 respectivamente).

\section{Análisis diferencial}
\label{sec:diferencial}

Para cada gen $g$ y cada cohorte $c$ se contrasta la diferencia de
expresión media entre los grupos mediante un test $t$ de Welch (que no
asume igual varianza entre grupos). Los $p$-valores se corrigen para
múltiples comparaciones con el método de Benjamini-Hochberg (FDR).

Se calcula además el tamaño de efecto mediante la $d$ de Cohen:
\begin{equation}
  d_g = \frac{\bar{x}_{g,\text{enfermo}} - \bar{x}_{g,\text{sano}}}{s_{p,g}}
\end{equation}
donde $s_{p,g}$ es la desviación típica combinada. Se considera un tamaño
de efecto \emph{grande} si $|d| > 0{,}8$, y \emph{muy grande} si
$|d| > 1{,}2$.

\section{Clasificadores supervisados}
\label{sec:clasificadores}

Los tres clasificadores empleados se implementan sobre \texttt{scikit-learn}
1.4 con las siguientes configuraciones base:

\begin{itemize}
  \item \textbf{Regresión logística L1:}
  \texttt{LogisticRegression(penalty='l1', solver='saga', C=0.5,
  max\_iter=5000, random\_state=42)}. La regularización se calibra en un
  pequeño estudio previo con búsqueda logarítmica en $C \in \{0{,}1;
  0{,}5; 1; 5\}$; el valor 0,5 ofrece el mejor compromiso entre número de
  genes seleccionados (mediana $\sim$300) y AUC en validación interna.
  \item \textbf{Random Forest:}
  \texttt{RandomForestClassifier(n\_estimators=300, max\_features='sqrt',
  random\_state=42)}.
  \item \textbf{Support Vector Machine:}
  \texttt{SVC(kernel='linear', C=1.0, probability=True, random\_state=42)}.
\end{itemize}

La tabla~\ref{tab:hiperparametros} resume los hiperparámetros y las
propiedades relevantes de cada clasificador.

\begin{table}[H]
  \centering
  \small
  \caption{Hiperparámetros y propiedades de los tres clasificadores.}
  \label{tab:hiperparametros}
  \begin{tabular}{lllll}
    \toprule
    \textbf{Modelo} & \textbf{Hiperparámetros} & \textbf{Regularización} & \textbf{Interpretabilidad} & \textbf{Determinismo} \\
    \midrule
    Regresión log. L1 & $C{=}0{,}5$, saga, $10^{4{,}7}$ iter. & $\ell_1$ (esparsa) & Alta (coefs.) & sí (seed 42) \\
    Random Forest & 300 árboles, $\sqrt{p}$ feat. & Implícita (bagging) & Media (importancias) & sí (seed 42) \\
    SVM lineal & $C{=}1{,}0$, kernel lineal & $\ell_2$ (densa) & Baja (pesos) & sí (seed 42) \\
    \midrule
    Baseline & Predice clase mayoritaria & — & — & determinista \\
    \bottomrule
  \end{tabular}
\end{table}

Como \emph{baseline} de comparación se emplea el clasificador
\emph{mayoritario informado}: siempre predice la clase más frecuente en el
conjunto de entrenamiento y su balanced accuracy es siempre 0,5. Esto
proporciona un umbral mínimo por debajo del cual el clasificador no aporta
información.

\section{Validación externa Leave-One-Dataset-Out}
\label{sec:lodo}

La validación se realiza mediante LODO. Para cada iteración $k$:

\begin{enumerate}
  \item Se selecciona la cohorte $c_k$ como test.
  \item Se entrena el modelo con todas las muestras de las cohortes
  restantes.
  \item Se evalúa en $c_k$ y se calculan AUC, balanced accuracy,
  sensibilidad, especificidad y ganancia respecto al baseline.
\end{enumerate}

Los resultados de las $k$ iteraciones se agregan como media (para las
cohortes evaluables) y se registran también individualmente en
\texttt{LODO\_HONESTO\_RESULTADOS.csv} y
\texttt{SUBTIPO\_LODO\_RESULTADOS.csv}. La figura~\ref{fig:lodo} ilustra
gráficamente el esquema de la validación LODO para un caso con cuatro
cohortes.

\begin{figure}[H]
  \centering
  \begin{tikzpicture}[node distance=0.35cm and 0.55cm]
    % Iteracion 1
    \node[paso, fill=orange!20, draw=orange!70!black] (t1) {Test\\C1};
    \node[paso, right=of t1] (a1) {Train C2};
    \node[paso, right=of a1] (a2) {Train C3};
    \node[paso, right=of a2] (a3) {Train C4};
    \node[right=of a3, font=\footnotesize] (r1) {\;$\Rightarrow$ AUC$_1$, BalAcc$_1$};
    % Iteracion 2
    \node[paso, below=of t1] (b1) {Train C1};
    \node[paso, fill=orange!20, draw=orange!70!black, right=of b1] (t2) {Test\\C2};
    \node[paso, right=of t2] (b2) {Train C3};
    \node[paso, right=of b2] (b3) {Train C4};
    \node[right=of b3, font=\footnotesize] (r2) {\;$\Rightarrow$ AUC$_2$, BalAcc$_2$};
    % Iteracion 3
    \node[paso, below=of b1] (c1) {Train C1};
    \node[paso, right=of c1] (c2) {Train C2};
    \node[paso, fill=orange!20, draw=orange!70!black, right=of c2] (t3) {Test\\C3};
    \node[paso, right=of t3] (c3) {Train C4};
    \node[right=of c3, font=\footnotesize] (r3) {\;$\Rightarrow$ AUC$_3$, BalAcc$_3$};
    % Iteracion 4
    \node[paso, below=of c1] (d1) {Train C1};
    \node[paso, right=of d1] (d2) {Train C2};
    \node[paso, right=of d2] (d3) {Train C3};
    \node[paso, fill=orange!20, draw=orange!70!black, right=of d3] (t4) {Test\\C4};
    \node[right=of t4, font=\footnotesize] (r4) {\;$\Rightarrow$ AUC$_4$, BalAcc$_4$};
    % Media
    \node[salida, below=0.7cm of d2, minimum width=8cm] (media)
      {Media LODO $=\frac{1}{k}\sum_{i=1}^{k}$ AUC$_i$, BalAcc$_i$};
    \draw[flecha] (r1.south east) to[out=-15, in=45] (media.north east);
    \draw[flecha] (r2.south east) to[out=-15, in=45] (media.east);
    \draw[flecha] (r3.south east) to[out=-15, in=25] (media.east);
    \draw[flecha] (r4.west) -- (media.east);
  \end{tikzpicture}
  \caption{Esquema de la validación Leave-One-Dataset-Out con $k=4$
  cohortes. En cada iteración, una cohorte (naranja) actúa como test y
  las demás como entrenamiento. Las métricas se agregan como media al
  final del proceso.}
  \label{fig:lodo}
\end{figure}

\section{Firma génica por consenso}
\label{sec:consenso}

La firma génica validada se construye por meta-análisis de consenso
multi-cohorte. Un gen entra en la firma si y sólo si cumple
simultáneamente en las tres cohortes evaluables:

\begin{enumerate}
  \item El signo del cambio (positivo o negativo) es idéntico.
  \item La magnitud del cambio, medida como $|d|$ de Cohen, es mayor que un
  umbral $\tau = 0{,}5$ (efecto medio) en todas las cohortes.
\end{enumerate}

El listado resultante se ordena por $d$ media entre cohortes y se
almacena en \texttt{FIRMA\_VALIDADA\_COMPLETA.csv}. Los 60 primeros genes
por ranking se exportan también a \texttt{FIRMA\_VALIDADA\_TOP60.csv} para
facilitar la exploración manual. La figura~\ref{fig:consenso} ilustra
visualmente el criterio de consenso sobre tres genes de ejemplo:

\begin{figure}[H]
  \centering
  \begin{tikzpicture}[node distance=0.45cm and 1cm]
    % Cabeceras
    \node[font=\small\bfseries] (h1) {Gen};
    \node[right=1.3cm of h1, font=\small\bfseries] (h2) {Cohorte A};
    \node[right=of h2, font=\small\bfseries]      (h3) {Cohorte B};
    \node[right=of h3, font=\small\bfseries]      (h4) {Cohorte C};
    \node[right=of h4, font=\small\bfseries]      (h5) {Decisión};

    % Gen 1: DSG3 - consenso
    \node[below=of h1, font=\small\ttfamily] (g1) {DSG3};
    \node[paso, right=1.3cm of g1, fill=green!15, minimum width=2cm, minimum height=0.7cm, font=\footnotesize] (a1) {$d = +4{,}12$};
    \node[paso, right=of a1, fill=green!15, minimum width=2cm, minimum height=0.7cm, font=\footnotesize] (b1) {$d = +4{,}05$};
    \node[paso, right=of b1, fill=green!15, minimum width=2cm, minimum height=0.7cm, font=\footnotesize] (c1) {$d = +4{,}36$};
    \node[salida, right=of c1, minimum width=2.4cm, minimum height=0.7cm, font=\footnotesize, fill=green!25, draw=green!70!black] (d1) {\checkmark\ Entra};

    % Gen 2: BRCA1 - signo opuesto
    \node[below=of g1, font=\small\ttfamily] (g2) {\emph{Gen X}};
    \node[paso, right=1.3cm of g2, fill=green!15, minimum width=2cm, minimum height=0.7cm, font=\footnotesize] (a2) {$d = +1{,}20$};
    \node[paso, right=of a2, fill=red!12, minimum width=2cm, minimum height=0.7cm, font=\footnotesize] (b2) {$d = -0{,}90$};
    \node[paso, right=of b2, fill=green!15, minimum width=2cm, minimum height=0.7cm, font=\footnotesize] (c2) {$d = +1{,}05$};
    \node[paso, right=of c2, minimum width=2.4cm, minimum height=0.7cm, font=\footnotesize, fill=red!12, draw=red!70!black] (d2) {$\times$ Signo no coincide};

    % Gen 3: magnitud insuficiente
    \node[below=of g2, font=\small\ttfamily] (g3) {\emph{Gen Y}};
    \node[paso, right=1.3cm of g3, fill=green!15, minimum width=2cm, minimum height=0.7cm, font=\footnotesize] (a3) {$d = +0{,}90$};
    \node[paso, right=of a3, fill=yellow!20, minimum width=2cm, minimum height=0.7cm, font=\footnotesize] (b3) {$d = +0{,}35$};
    \node[paso, right=of b3, fill=green!15, minimum width=2cm, minimum height=0.7cm, font=\footnotesize] (c3) {$d = +1{,}10$};
    \node[paso, right=of c3, minimum width=2.4cm, minimum height=0.7cm, font=\footnotesize, fill=yellow!30, draw=orange!80!black] (d3) {$\times$ $|d| < 0{,}5$ en B};

    % Regla al pie
    \node[grupo, fit=(g1) (d3), inner sep=6pt] (box) {};
    \node[below=0.4cm of box, font=\footnotesize\itshape, text width=13cm, align=center]
      {Regla de consenso: un gen entra en la firma sólo si mantiene el
      mismo signo y $|d| > 0{,}5$ en las 3 cohortes evaluables.};
  \end{tikzpicture}
  \caption{Regla del meta-análisis por consenso. Sobre tres genes de
  ejemplo, sólo \gene{DSG3} pasa el filtro: mantiene signo positivo y
  $|d|$ grande en las tres cohortes. \emph{Gen X} tiene signo opuesto
  en la cohorte B; \emph{Gen Y} tiene signo consistente pero magnitud
  insuficiente en B.}
  \label{fig:consenso}
\end{figure}

\section{Panel mínimo}
\label{sec:panel_minimo}

Para determinar el panel mínimo se realiza un análisis de sensibilidad:
para cada tamaño $k \in \{3, 5, 8, 10, 15, 20, 30, 50, 100, 250, 500,
1174\}$, se selecciona el top-$k$ de la firma validada y se evalúa
mediante LODO. El panel mínimo se define como el menor $k$ cuya AUC media
está a menos de 0,01 del máximo alcanzado por la curva completa.

La curva de AUC en función de $k$ se almacena en
\texttt{PANEL\_MINIMO\_CURVA.csv} y su resumen (panel mínimo, AUC del
panel, AUC de la firma completa) se serializa en el JSON
\texttt{FIRMA\_VALIDADA\_RESUMEN.json}.

\section{Validación contra IHC clínica}
\label{sec:validacion_ihc}

El panel diagnóstico recomendado por la OMS para NSCLC~\cite{travis2015who}
se codifica en el pipeline como una lista fija de 20 genes:

\begin{itemize}
  \item \textbf{Marcadores escamosos} (12): \gene{KRT5}, \gene{KRT6A},
  \gene{KRT6B}, \gene{KRT13}, \gene{KRT14}, \gene{TP63}, \gene{DSG3},
  \gene{DSC3}, \gene{SOX2}, \gene{PKP1}, \gene{CALML3}, \gene{S100A2}.
  \item \textbf{Marcadores adenocarcinoma} (8): \gene{NAPSA},
  \gene{NKX2-1}, \gene{SFTPB}, \gene{SFTPC}, \gene{SFTPA1}, \gene{SLC34A2},
  \gene{MUC1}, \gene{CEACAM6}.
\end{itemize}

La validación consiste en intersectar esta lista con
\texttt{FIRMA\_VALIDADA\_COMPLETA.csv} y contar cuántos marcadores se han
recuperado sin declararlos al framework durante el entrenamiento.

\section{Interfaz web y asistente}
\label{sec:interfaz}

La interfaz web se implementa en Streamlit 1.35 y se estructura como una
aplicación multipágina con dos vistas:

\begin{itemize}
  \item \textbf{Portada (\texttt{/})}: hoja de presentación con el título
  del trabajo y una tarjeta resumen de cifras clave.
  \item \textbf{Framework (\texttt{/framework})}: dashboard con siete
  capítulos (Inicio, Asistente, Introducción, Metodología, Resultados,
  Cohortes, Conclusiones), cada uno reproduciendo la estructura de la
  memoria escrita y leyendo las cifras dinámicamente de los CSV/JSON del
  pipeline.
\end{itemize}

El asistente conversacional integra el mismo modelo Llama~3.3-70b
empleado en la curación clínica, alimentado con un prompt sistema
(\texttt{contexto\_tfm.py}) que contiene el contexto de resultados leído
dinámicamente y una serie de reglas explícitas para acotar el
comportamiento: fidelidad a los datos, prohibición de inventar cifras,
enfoque en resultados biológicos y de rendimiento, y encuadre de los
criterios de inclusión como metodología y no como resultado principal.

\section{Reproducibilidad}
\label{sec:reproducibilidad}

Todo el análisis puede reproducirse desde la raíz del repositorio con la
siguiente secuencia de comandos:

\begin{lstlisting}[style=bash, caption={Reproducción completa del pipeline.}]
python agentes/paso1_descarga.py
python agentes/paso2_metadata.py
python agentes/paso3_diferencial.py
python agentes/paso4_informe.py
python agentes/paso5_graficos.py
python agentes/paso6_dashboard.py
python agentes/paso7_orquestador.py
python agentes/paso13_subtipo_lodo.py
python agentes/paso14_auditoria_datos.py
python agentes/paso15_lodo_honesto.py
python agentes/paso16_composicion_vs_biologia.py
python agentes/paso17_falacia_folds.py
python agentes/paso18_subtipo_casos_dificiles.py
python agentes/paso19_firma_validada.py
python agentes/generar_figuras_auditoria.py
streamlit run agentes/paso12_web_chatbot.py
\end{lstlisting}

La única dependencia externa es la disponibilidad de una clave API de Groq
para la curación LLM y para el asistente conversacional. Sin esta clave el
pipeline puede reproducirse siempre que exista una versión previa del
fichero \texttt{AUDITORIA\_COHORTES.csv} con las etiquetas asignadas, ya
que los pasos posteriores no vuelven a consultar el LLM.
